% Created 2026-05-04 lun 16:50
% Intended LaTeX compiler: pdflatex
\documentclass[10pt,a4paper]{report}
\usepackage[utf8]{inputenc}
\usepackage{lmodern}
\usepackage[T1]{fontenc}
\usepackage[top=1in, bottom=1.in, left=1in, right=1in]{geometry}
\usepackage{graphicx}
\usepackage{longtable}
\usepackage{float}
\usepackage{wrapfig}
\usepackage{rotating}
\usepackage[normalem]{ulem}
\usepackage{amsmath}
\usepackage{textcomp}
\usepackage{marvosym}
\usepackage{wasysym}
\usepackage{amssymb}
\usepackage{amsmath}
\usepackage[theorems, skins]{tcolorbox}
\usepackage[version=3]{mhchem}
\usepackage[numbers,super,sort&compress]{natbib}
\usepackage{natmove}
\usepackage{url}
\usepackage[cache=false]{minted}
\usepackage[strings]{underscore}
\usepackage[linktocpage,pdfstartview=FitH,colorlinks,
linkcolor=blue,anchorcolor=blue,
citecolor=blue,filecolor=blue,menucolor=blue,urlcolor=blue]{hyperref}
\usepackage{attachfile}
\usepackage{setspace}
\author{Marcos Bujosa}
\date{\today}
\title{calcprop — cálculo proposicional}
\begin{document}

\tableofcontents

\part{Aparato lógico}
\label{sec:org93fd436}

\chapter{Estructura del paquete \texttt{calcprop}}
\label{sec:org92b2248}

\begin{minted}[frame=lines,fontsize=\scriptsize,linenos=]{python}
<<Aparato lógico de la librería>>
\end{minted}

\begin{minted}[frame=lines,fontsize=\scriptsize,linenos=]{python}
<<Definición de la clase Proposicion>>
<<Definición de la subclase v>>
<<Definición de alguno y unoDe>>
<<Definición de noDerivable>>
<<Definición de la función span>>
<<Definición de la función extension>>
<<Definición de la función refuta>>
<<Definición de la función test>>
\end{minted}

El \texttt{<{}<{}Aparato auxiliar para la creación de problemas de opción múltiple>{}>{}}
se desarrolla en el paquete \texttt{calcprop-quiz}.
\chapter{Clase \texttt{Proposicion} y su subclase \texttt{v} (proposición atómica)}
\label{sec:orga6c06dc}

Las proposiciones son expresiones formadas con números, strings y los símbolos
\begin{center}
\texttt{\&},\quad \texttt{|},\quad \texttt{-},\quad  \texttt{v( )},\quad \texttt{>{}>{}},\quad \texttt{**},\quad \texttt{alguno}\quad y\quad  \texttt{unoDe}
\end{center}
construidas según las siguientes reglas:
\begin{enumerate}
\item Si \texttt{D} es un número o un string (cadena de caracteres) entonces \texttt{v(D)} es una
subclase de \texttt{Proposicion} que denominaremos \emph{``variable proposicional''} o \emph{``proposición atómica''}.
\item Si \texttt{P}, \texttt{Q} y \texttt{R} son proposiciones también lo son:
\begin{description}
\item[{\texttt{P\&Q}}] (conjunción)
\item[{\texttt{P|Q}}] (disyunción)
\item[{\texttt{-P}}] (negación)
\item[{\texttt{P>{}>{}Q}}] (implicación)
\item[{\texttt{P**Q}}] (equivalencia, es decir: \emph{si y solo si}).
\end{description}
Para facilitar la elaboración expresiones, también incluimos dos operadores lógico menos habituales:
\begin{description}
\item[{\texttt{alguno(P,Q,R)}}] significa que alguna de las proposiciones de la lista es verdadera y
\item[{\texttt{unoDe(P,Q,R)}}] significa que solo una de las proposiciones de la lista es verdadera.
\end{description}
\end{enumerate}

Ejemplo de proposición: ``llueve o no llueve''
\begin{minted}[frame=lines,fontsize=\scriptsize,linenos=]{python}
v('llueve') | -v('llueve')
\end{minted}

\phantomsection
\label{}
\begin{verbatim}
v('llueve') | -v('llueve')
\end{verbatim}


Necesitamos que las proposiciones del tipo \texttt{v(D)} sean \href{https://docs.python.org/3/glossary.html\#term-hashable}{\emph{``hash-eables''}}, para poder ser utilizadas como claves en un diccionario.
Para ello necesitamos los métodos \texttt{\_\_hash\_\_} y \texttt{\_\_eq\_\_}.

Definamos la clase \texttt{Proposicion}
\begin{minted}[frame=lines,fontsize=\scriptsize,linenos=]{python}
class Proposicion:
    def __init__(self,data):
        self.data = data

    def __hash__(self):
        return hash(self.data)

    def __eq__(self, another):
        return hasattr(another, 'data') and self.data == another.data

    <<Operadores lógicos de la clase Proposicion>>
    <<Método de representación de la clase Proposicion>>
\end{minted}
\begin{enumerate}
\item Las operaciones: conjunción, disyunción, negación, implicación y equivalencia.
\label{sec:org0b17f94}

Definamos como métodos dentro de la clase \texttt{Proposicion} los cinco operadores lógicos habituales.
Primero aquellos que usan \emph{``métodos mágicos''}, es decir, aquellos que tiene asociado un símbolo para llamar a las operaciones:
\begin{minted}[frame=lines,fontsize=\scriptsize,linenos=]{python}
def __and__(self,other):
    return Proposicion(['and', self, other])
def __or__(self,other):
    return Proposicion(['or', self, other])
def __neg__(self):
    return Proposicion(['not', self])
def __rshift__(self,other):
    return Proposicion(['implica', self, other])
def __pow__(self,other):
    return Proposicion(['equivale', self, other])

\end{minted}
\end{enumerate}
\item Las operaciones: \texttt{unoDe} y \texttt{alguno}.
\label{sec:orgf26e2ea}

Para las dos siguientes operaciones\footnote{no habituales, aunque convenientes pues facilitan la expresión de algunas proposiciones} no usamos ningún \emph{``método mágico''}.
Así que definirlas dentro de la clase nos obligaría a una incómoda escritura del tipo: \texttt{Proposicion.alguno(P1,P2,P3)} o \texttt{Proposicion.unoDe(P1,P2,P3)}.
Para evitarlo, las definimos fuera de la clase; así podremos escribir sencillamente \texttt{alguno(P1,P2,P3)} y \texttt{unoDe(P1,P2,P3)}.
\begin{minted}[frame=lines,fontsize=\scriptsize,linenos=]{python}
def alguno(X,*args):
    return Proposicion(['alguno', X] + [a for a in args])
def unoDe(X,*args):
    return Proposicion(['unoDe',  X] + [a for a in args])

\end{minted}
\item La operación: \texttt{noDerivable}.
\label{sec:org939fa92}

Esta función no es una operación lógica (como si lo son las anteriores).
Es un ``añadido'' cuya motivación es la siguiente:

El esquema típico de una pregunta es: el enunciado asume que las premisas \texttt{E1}, \texttt{E2},\dots{} \texttt{En} son ciertas.
Y se puede puede preguntar si \texttt{P} es consecuencia lógica de las premisas.
Si lo es, \texttt{P} es \texttt{True} (es decir, \(\texttt{premisas}\Rightarrow\texttt{P}\)).
Negar esto no es decir que \texttt{-P} es consecuencia de las premisas, sino que \texttt{P} no es derivable de las premisas (es decir, \(\texttt{premisas}\not\Rightarrow\texttt{P}\)).
Así que añadimos la operación \texttt{noDerivable}.
Esto nos permite, por ejemplo, que aunque para algunos enunciados se pueda preguntar si una matriz es diagonalizable, para otros evitemos la pregunta si no es deducible del enunciado.

\begin{minted}[frame=lines,fontsize=\scriptsize,linenos=]{python}
def noDerivable(X):
    return Proposicion(['noDerivable', X])

\end{minted}
\end{enumerate}
\section{La subclase \texttt{v} para proposiciones atómicas (i.e., para variables proposicionales)}
\label{sec:org0e15e5c}

\begin{minted}[frame=lines,fontsize=\scriptsize,linenos=]{python}
class v(Proposicion):
    def __init__(self, data):
        super().__init__(data)

\end{minted}
\chapter{Valoración de proposiciones}
\label{sec:org21447d4}

Una valoración es una función \texttt{T} que va del conjunto de proposiciones al conjunto \texttt{\{True,False\}} (es decir, a cada proposición le asigna el valor verdadero o el valor falso) y que verifica las siguientes propiedades: si \texttt{P} y \texttt{Q} son proposiciones entonces

\begin{enumerate}
\item \texttt{T(P \& Q) = True}\quad únicamente en el caso de que\quad \texttt{T(P) = True}\quad y\quad \texttt{T(Q) = True}.
\item \texttt{T(P | Q) = False}\quad únicamente en el caso de que\quad \texttt{T(P) = False}\quad y\quad \texttt{T(Q) = False}.
\item \texttt{T(-P) = True}\quad únicamente en el caso de que\quad \texttt{T(P) = False}.
\end{enumerate}
Por tanto:
\fbox{\texttt{T(P \& Q) = T(P) \& T(Q)}},\qquad
\fbox{\texttt{T(P | Q) = T(P) | T(Q)}}\quad y \quad
\fbox{\texttt{T(-P) = not T(P)}}.
\begin{enumerate}
\item \texttt{T(P >{}>{} Q) = False}\quad únicamente en el caso en que\quad \texttt{T(P) = True}\quad y\quad \texttt{T(Q) = False}.
\item \texttt{T(P ** Q) = True}\quad únicamente en el caso en que\quad \texttt{T(P) = T(Q)}.
\end{enumerate}
Por tanto\quad
\fbox{\texttt{T(P >> Q) = T(-P | Q)}}
\quad y\quad
\fbox{\texttt{T(P ** Q) = T( (P>>Q) \& (Q>>P) )}}.

Y si \texttt{A1,...An} son \texttt{n} proposiciones, entonces
\begin{enumerate}
\item \texttt{alguno(A1,...,An)} es \texttt{True}\quad si y solo si\quad \texttt{A1 | alguno(A2,...,An)} es \texttt{True},
\item \texttt{unoDe(A1,...,An)} es \texttt{True}\quad si y solo si\quad \texttt{(A1 \& -alguno(A2,...,An) | (-A1 \& unoDe(A2,...,An))} es \texttt{True};
\end{enumerate}
Por tanto
\begin{center}
\fbox{\texttt{T( alguno(A1,...,An) ) = T( A1 | alguno(A2,...,An) )}}
\end{center}
y
\begin{center}
\fbox{\texttt{T( unoDe(A1,...,An) ) = T( (A1 \& -alguno(A2,...,An)) | (-A1 \& unoDe(A2,...,An)) )}}.
\end{center}
Cualquier función \texttt{F} que vaya del conjunto de proposicionales atómicas (es decir, del tipo \texttt{v('algo')}) al conjunto \texttt{\{True,False\}}
se puede extender al resto proposiciones\footnote{En la función \texttt{span} de Python hemos elegido el criterio de que a toda proposición \texttt{Prop} no incluida explícitamente en el diccionario \texttt{F} le asignamos el valor \texttt{False} para así completar el dominio de las proposiciones atómicas.
Esta decisión ha sido completamente arbitraria, y se ha incluido para que efectivamente el dominio de \texttt{span} incluya todas las proposiciones atómicas (i.e., variables proposicionales).
En cualquier caso \texttt{span} se ha programado a modo de ilustración, pues no la emplearemos para nada.}
formando una valoración \texttt{T=span(F)}, aplicando las reglas de arriba.
Dicha extensión es única.

Vamos a implementar como funciones los diccionarios de Python, pues los diccionarios son una colección de pares \texttt{'Etiqueta':valor} donde a cada etiqueta solo se le asigna un valor único.
Por ejemplo, si definimos el siguiente diccionario \texttt{\{'a':1 , 'a':0 , 'b':2}\}, obtenemos un diccionario con dos elementos, donde a \texttt{'a'} se le a asignado solo uno de los dos valores introducidos (en concreto el último).
\begin{minted}[frame=lines,fontsize=\scriptsize,linenos=]{python}
print( {'a':1 , 'a':0 , 'b':2} )
\end{minted}

\phantomsection
\label{}
\begin{verbatim}
{'a': 0, 'b': 2}
\end{verbatim}
\section{Función \texttt{span}}
\label{sec:org3073dd3}

Función \texttt{span(F, Prop)} indica si \texttt{Prop} es \texttt{True} o es \texttt{False} a partir de la función parcial \texttt{F}.
\begin{itemize}
\item \texttt{F}: es un diccionario (función parcial) que asigna valores \texttt{True} o \texttt{False} a varias proposiciones atómicas.
\item \texttt{Prop}: es una \texttt{Proposicion} .
\end{itemize}
Así obtenemos los siguientes resultados lógicos:
\begin{minted}[frame=lines,fontsize=\scriptsize,linenos=]{python}
F = { v('llueve'):True }
span(F, v('llueve') >> -v('llueve') )
\end{minted}

\phantomsection
\label{}
\begin{verbatim}
False
\end{verbatim}

y por otra parte
\begin{minted}[frame=lines,fontsize=\scriptsize,linenos=]{python}
span(F, -v('llueve') >> v('llueve') )
\end{minted}

\phantomsection
\label{}
\begin{verbatim}
True
\end{verbatim}
\subsection{Implementación}
\label{sec:orgb63986a}

Si una proposición \texttt{Prop} (de tipo atómico) está en la función parcial \texttt{F} y además es verdadera, \texttt{span} devuelve \texttt{True}.
En caso contrario devuelve \texttt{False} (tanto si \texttt{Prop} no está en \texttt{F} como si lo está pero es \texttt{False}).
Si \texttt{Prop} no es de tipo atómico, entonces se aplican las reglas para reducir \texttt{Prop} a proposiciones de tamaño más pequeño.

\begin{minted}[frame=lines,fontsize=\scriptsize,linenos=]{python}
def span(F, Prop):
    if type(Prop) == type(v('')):
        return (Prop in F) and F[Prop]
    <<Aplicación de las reglas para reducir el tamaño de Prop>>

\end{minted}

\begin{minted}[frame=lines,fontsize=\scriptsize,linenos=]{python}
elif Prop.data[0] == 'and':
    return span(F, Prop.data[1]) and span(F, Prop.data[2])

elif Prop.data[0] == 'or':
    return span(F, Prop.data[1]) or span(F, Prop.data[2])

elif Prop.data[0] == 'not':
    return not span(F, Prop.data[1])

elif Prop.data[0] == 'implica':
    return (not span(F, Prop.data[1])) or span(F, Prop.data[2])

elif Prop.data[0] == 'equivale':
    return span(F, Prop.data[1]) == span(F, Prop.data[2])

elif Prop.data[0] == 'alguno':
    return span(F,Prop.data[1]) or span(F,alguno(Prop.data[2:])) if len(Prop.data)>2\
      else span(F,Prop.data[1])

elif Prop.data[0] == 'unoDe':
    A = Prop.data[1]
    if len(Prop.data)==2:
        return span(F, A)
    else:
        B = Proposicion(['alguno'] + Prop.data[2:])
        C = Proposicion(['unoDe']  + Prop.data[2:])
        return span(F, (A&-B) | (-A&C) )
\end{minted}
\chapter{Función  \texttt{extension} que extiende una función parcial}
\label{sec:org707aa3c}

Ahora queremos definir una función de Python (que denominamos \texttt{extension}) que,
dadas una lista \texttt{objetivo} y un diccionario con \texttt{valores},
para una colección de proposiciones atómicas haga crecer (extienda) el diccionario \texttt{valores} con nuevas proposiciones atómicas, y sus correspondientes valores,
necesarios para que se cumplan los pares contenidos en la lista \texttt{objetivo};
o bien que devuelva el diccionario vacío \texttt{\{\}} en caso de que sea imposible hacer cumplir dicha lista \texttt{objetivo}.

La diferencia entre la lista \texttt{objetivo} y los \texttt{valores} es que en \texttt{objetivo} se asignan valores a proposiciones (arbitrarias) pero en \texttt{valores} solo se asignan valores a \emph{proposiciones atómicas}.

El diccionario extendido es un \emph{modelo lógico} en el que se cumplen la lista \texttt{objetivo} y la función parcial \texttt{valores}.
Ello no quiere decir que no puedan existir otros modelos que consigan lo mismo.

Ejemplo:
\begin{minted}[frame=lines,fontsize=\scriptsize,linenos=]{python}
print( extension( [ (v('llueve') >> -v('llueve') , True) ], {v('basilisco') : True}) )
\end{minted}

\phantomsection
\label{}
\begin{verbatim}
{v('basilisco'): True, v('llueve'): False}
\end{verbatim}



Ejemplo:
\begin{minted}[frame=lines,fontsize=\scriptsize,linenos=]{python}
print( extension( [ (v('llueve'),True), (-v('llueve'),True)], {v('basilisco') : True}) )
\end{minted}

\phantomsection
\label{}
\begin{verbatim}
{}
\end{verbatim}
\section{Implementación de la función  \texttt{extension}}
\label{sec:org131622c}

Función \texttt{extension(objetivo, valores, premisas=[])}

\begin{itemize}
\item \texttt{objetivo}: es una lista de pares \texttt{(P,VF)} donde \texttt{P} es una proposición y \texttt{VF} es \texttt{True} o \texttt{False}

\item \texttt{valores}: es un diccionario (función parcial) que asigna valores \texttt{True} o \texttt{False} a variables proposicionales

\item \texttt{premisas}: es la lista de premisas a partir de la cual se deduce si una proposición es derivable o no (es decir, es solo para usar la operación \texttt{noDerivable})
\end{itemize}

La función \texttt{extension} tiene dos posibles salidas:
o bien devuelve un modelo lógico \texttt{M} que extiende a \texttt{valores} de tal forma que \texttt{span(M,P)=VF} para cada par \texttt{(P,VF)} de \texttt{objetivo},
o el diccionario vacío \texttt{\{\}} si no existe tal modelo \texttt{M}.
Es decir,
busca una forma de dar valores a las proposiciones atómicas de modo que todas las proposiciones contenidas en el \texttt{objetivo} alcancen el valor deseado.
\subsection{Implementación}
\label{sec:org1d8a1d3}

Si la lista \texttt{objetivo} no tiene pares, con la propia función \texttt{valores}, hemos acabado.

La definición de la función \texttt{extension} es recursiva.
La mecánica es la siguiente: se sustituye el proposición del primer objetivo por proposiciones de menor tamaño\footnote{Una lista de símbolos más corta}, y entonces se llama a si misma.
Este procedimiento reduce las proposiciones hasta obtener proposiciones atómicas.

Así pues, llamamos \texttt{Obj} al primer objetivo de la lista, \texttt{P} a la correspondiente proposición y \texttt{VF} a su valor objetivo.
Por tanto \texttt{Obj = (P,VF)} es el primer par de la lista \texttt{objetivo}.

Comprobamos que \texttt{P.data} es una \texttt{list} y tratamos los distintos casos de reducción de predicado.

\begin{minted}[frame=lines,fontsize=\scriptsize,linenos=]{python}
def extension(objetivo, valores, premisas=[]):
    if objetivo == []:
        return valores

    Obj = objetivo[0]
    P   = Obj[0]
    VF  = Obj[1]
    if isinstance(P.data,list):
        <<Reducción de una proposición de tipo "not">>
        <<Reducción de una proposición de tipo "and">>
        <<Reducción de una proposición de tipo "or">>
        <<Reducción de una proposición de tipo "implica">>
        <<Reducción de una proposición de tipo "equivale">>
        <<Reducción de una proposición de tipo "alguno">>
        <<Reducción de una proposición de tipo "unoDe">>
        <<Proposición de tipo "noDerivable">>
    <<Proposición atómica v>>
\end{minted}
\section{Reducción de proposiciones NO atómicas}
\label{sec:org039e846}

\begin{enumerate}
\item Reducción de una proposición de tipo ``not''.
\label{sec:org2771e45}

En el caso de que \texttt{Obj = (-A,True)}, nos vale la misma función que obtenemos si remplazamos \texttt{Obj} por \texttt{(A,False)}.
Y en el caso de que \texttt{Obj = (-A,False)}, nos vale la misma función si remplazamos \texttt{Obj} por \texttt{(A,True)}.
\begin{minted}[frame=lines,fontsize=\scriptsize,linenos=]{python}
if P.data[0] == 'not' :
    A = P.data[1]
    return extension( [(A, not VF)] + objetivo[1:], valores, premisas )
\end{minted}
\item Reducción de una proposición de tipo ``and''.
\label{sec:org3c7622f}

En el caso de que \texttt{Obj=(A \& B, True)}, nos vale la misma función que obtenemos si remplazamos \texttt{Obj} por \texttt{(A,True), (B,True)}.
Y en el caso de que \texttt{Obj=(A \& B, False)}, vale tanto la función que obtenemos de remplazar \texttt{Obj} por \texttt{(A,False)} como la que obtenemos al remplazar \texttt{Obj} por \texttt{(B,False)}.

Si la primera alternativa resulta en una extensión vacía se devuelve la segunda.
\begin{minted}[frame=lines,fontsize=\scriptsize,linenos=]{python}
elif P.data[0] == 'and':
    A = P.data[1]
    B = P.data[2]
    if(VF):
        return extension([ (A,True), (B,True) ] + objetivo[1:], valores, premisas)
    else:
        r = extension( [(A,False)] + objetivo[1:],valores,premisas)
        return extension([(B,False)]+objetivo[1:],valores,premisas) if not r else r

\end{minted}
\item Reducción de una proposición de tipo ``or''.
\label{sec:org854e9da}

En el caso de que \texttt{Obj=( A|B, False)}, nos vale la misma función que obtenemos si remplazamos \texttt{Obj} por \texttt{(A,False), (B,False)}.
Y en el caso de que \texttt{Obj=( A\&B, True)}, vale tanto la función que obtenemos de remplazar \texttt{Obj} por \texttt{(A,True)} como la que obtenemos al remplazar \texttt{Obj} por \texttt{(B,True)}.

Si la primera alternativa resulta en una extensión vacía se devuelve la segunda.
\begin{minted}[frame=lines,fontsize=\scriptsize,linenos=]{python}
elif P.data[0] == "or":
    A = P.data[1]
    B = P.data[2]
    if not VF:
        return extension([(A,False),(B,False)] + objetivo[1:], valores, premisas)
    else:
        r = extension([(A,True)] + objetivo[1:], valores, premisas)
        return extension([(B,True)]+objetivo[1:],valores, premisas) if not r else r

\end{minted}
\item Reducción de una proposición de tipo ``implica''.
\label{sec:orgf054ee0}

En el caso de que \texttt{Obj=( A>{}>{}B, VF)}, nos vale la misma función que obtenemos si emplazamos \texttt{Obj} por \texttt{( -A|B, VF)}.
\begin{minted}[frame=lines,fontsize=\scriptsize,linenos=]{python}
elif P.data[0] == "implica":
    A = P.data[1]
    B = P.data[2]
    return extension([ ( -A|B, VF) ] + objetivo[1:], valores, premisas)
\end{minted}
\item Reducción de una proposición de tipo ``equivale''.
\label{sec:org54e8a53}

En el caso de que \texttt{Obj=( A**B, VF)}, nos vale la misma función que obtenemos si remplazamos \texttt{Obj} por \texttt{( (A>{}>{}B) \& (B>{}>{}A), VF)}.
\begin{minted}[frame=lines,fontsize=\scriptsize,linenos=]{python}
elif P.data[0] == "equivale":
    A = P.data[1]
    B = P.data[2]
    return extension([( (A>>B) & (B>>A), VF)] + objetivo[1:], valores, premisas)
\end{minted}
\item Reducción de una proposición de tipo ``alguno''.
\label{sec:org2ddb761}

En caso de que \texttt{Obj = (alguno(A), VF)} (i.e., la lista de alternativas solo contiene la proposición \texttt{A}) nos vale la misma función que \texttt{(A, VF)}.
Si por el contrario \texttt{Obj = (alguno(A1,...An) ,VF)}, con \texttt{n>1}, nos vale la misma función que obtenemos si remplazamos \texttt{Obj} por \texttt{( A1|alguno(A2,...,An), VF)} (i.e., la primera o alguna de las demás).
\begin{minted}[frame=lines,fontsize=\scriptsize,linenos=]{python}
elif P.data[0] == "alguno":
    A = P.data[1]
    if len(P.data)==2:
        return extension([(A,VF)] + objetivo[1:], valores, premisas)
    else:
        B = Proposicion(["alguno"] + P.data[2:])
        return extension([( A|B, VF)] + objetivo[1:], valores, premisas)
\end{minted}
\item Reducción de una proposición de tipo ``unoDe''.
\label{sec:orgffdc323}

En caso de que \texttt{Obj=(unoDe(A), VF)} (i.e., la lista de alternativas solo contiene la proposición \texttt{A}) nos vale la misma función que \texttt{(A,VF)}.
Si por el contrario \texttt{Obj=(unoDe(A1,...An), VF)}, con \texttt{n>1}, nos vale la misma función que obtenemos si remplazamos \texttt{Obj} por \texttt{( A1 \& -alguno(A2,...An) | -A1 \& unoDe(A2,...,An), VF)}.
\begin{minted}[frame=lines,fontsize=\scriptsize,linenos=]{python}
elif P.data[0] == "unoDe":
    A = P.data[1]
    if len(P.data)==2:
        return extension([(A,VF)] + objetivo[1:], valores, premisas)
    else:
        B = Proposicion(["alguno"] + P.data[2:])
        C = Proposicion(["unoDe"] + P.data[2:])
        return extension([( (A&-B) | (-A&C), VF)] + objetivo[1:], valores, premisas)
\end{minted}
\end{enumerate}
\section{Proposición no derivable de un conjunto de \texttt{premisas}}
\label{sec:org7f8fd66}

La función \texttt{extension} tiene un tercer argumento (\texttt{premisas}) que es un ``añadido'' para decidir si una proposición es \texttt{noDerivable} de la lista de \texttt{premisas}.

Derivable significa que la proposición es consecuencia lógica de las \texttt{premisas}, por consiguiente es imposible que sea falso si las \texttt{premisas} son ciertas.

\begin{minted}[frame=lines,fontsize=\scriptsize,linenos=]{python}
elif P.data[0] == "noDerivable":
    A = P.data[1]
    valoresExt = valores.copy()
    valoresExt[v(repr(P))] = VF

    derivable = not extension([(A,False)] + premisas, {}, premisas)

    return extension(objetivo[1:],valores2,premisas) if not derivable == VF else {}

\end{minted}
\section{\ldots{} cuando llegamos a una proposición atómica (o \emph{variable proposicional})}
\label{sec:org21aa496}

En el caso de que \texttt{Obj = (v(D), VF)}, es cuando el algoritmo recursivo pasará a intentar un nuevo objetivo (extendiendo el diccionario \texttt{valores} si fuera necesario), o bien parará por ser el par objetivo \texttt{(v(D), VF)} imposible.

Primero se comprueba si el par \texttt{(v(D), VF)} ya estaba en el diccionario \texttt{valores}, es decir, si \texttt{valores} contiene la proposición atómica \texttt{v(D)} y con el valor \texttt{VF}.
En tal caso, se ``salta'' al siguiente objetivo (es decir se llama al \texttt{extension} con el mismo diccionario \texttt{valores} pero quitando de la lista \texttt{objetivo} el primer elemento.
Si \texttt{v(D)} pertenece al diccionario \texttt{valores} pero \textbf{con un valor distinto} a \texttt{VF}, quiere decir que el objetivo es imposible de alcanzar, por lo que se devuelve un diccionario vacío, y la recursión ha terminado.

Si \texttt{P} es una proposición atómica que no estaba en el diccionario \texttt{valores}, entonces creamos un nuevo diccionario extendido (\texttt{valoresExt}) que incluya \texttt{P} y su valor \texttt{VF}; y entonces llamamos a \texttt{extension} pero quitando el primer objetivo de la lista y usando un diccionario ``extendido'' (\texttt{valoresExt}).
Para evitar los efectos colaterales de modificar un diccionario, modificamos una copia

\begin{minted}[frame=lines,fontsize=\scriptsize,linenos=]{python}
else:
    if P in valores:
        return extension(objetivo[1:], valores, premisas) if valores[P]==VF else {}
    else:
        valoresExt    = valores.copy()
        valoresExt[P] = VF
        return extension(objetivo[1:], valoresExt, premisas)
\end{minted}
\chapter{Función que \texttt{refuta} una proposición a partir de un conjunto de \texttt{premisas}}
\label{sec:orgd9997d9}

Refutar una proposición, consiste en asignar un valor a cada variable proposicional (cada proposición atómica), de tal modo que el resultado sea falso.

La función \texttt{refuta} construye un modelo \texttt{M} (asignación de valores a las proposiciones atómicas) de forma que \texttt{span(M,P)=False} y \texttt{span(M,h)=True} para todo \texttt{h} en \texttt{premisas}.

Si esto no es posible refutar \texttt{P} a partir de las \texttt{premisas}, entonces la función \texttt{refuta} devuelve un modelo (un diccionario) vacío; esto quiere decir que la proposición \texttt{P} es irrefutable, es decir, \(\text{=premisas=}\Rightarrow\text{=P=}\).
Es decir, si se refuta es que no es derivable de las premisas (no es consecuencia lógica de las premisas).
\begin{minted}[frame=lines,fontsize=\scriptsize,linenos=]{python}
def refuta(P, premisas=[]):
    h=[(Q,True) for Q in premisas]
    return extension([(P,False)]+h, {}, h)

\end{minted}
\chapter{Función  \texttt{test}}
\label{sec:org53e521d}

Función \texttt{test(P,premisas)}

\begin{itemize}
\item \texttt{P}: es una proposición o cualquiera de los valores \texttt{True}, \texttt{False}.
\item \texttt{premisas}: es una lista de proposiciones
\end{itemize}

La función \texttt{test} se utiliza tanto para saber si se da una precondición, como para saber el resultado correcto de una cuestión.

En el caso de que \texttt{P} sea una proposición, sólo devuelve \texttt{True} si es imposible refutarlo a partir de las \texttt{premisas}, es decir, es imposible que \texttt{P} sea falso si las \texttt{premisas} son verdad (\texttt{premisas} \(\Rightarrow\) \texttt{P}, i.e., \texttt{P} es derivable).

\begin{minted}[frame=lines,fontsize=\scriptsize,linenos=]{python}
def test(P,premisas=[]):
    if P == True:
        return True
    if P == False:
        return False
    else:
        return True if not refuta(P, premisas) else False

\end{minted}
\chapter{Representación de la clase \texttt{Proposicion}}
\label{sec:org09147d7}

\begin{minted}[frame=lines,fontsize=\scriptsize,linenos=]{python}
def __repr__(self):
    if isinstance(self.data,list):
        if self.data[0] == "and":
            return repr(self.data[1]) + " & "+ repr(self.data[2])
        if self.data[0] == "or":
            return repr(self.data[1]) + " | "+ repr(self.data[2])
        if self.data[0] == "not":
            return "-" + repr(self.data[1])
        if self.data[0] == "implica":
            return repr(self.data[1]) + " >> "+ repr(self.data[2])
        if self.data[0] == "equivale":
            return repr(self.data[1]) + " ** "+ repr(self.data[2])
        if self.data[0] == "unoDe":
            return "unoDe("+','.join([repr(x) for x in self.data[1:]]) +")"
        if self.data[0] == "noDerivable":
            return "noDerivable("+repr(self.data[1]) +")"

    else:
        return 'v('+repr(self.data)+')'
\end{minted}
\end{document}
